\section{The secondary cusp on the whole interval}\label{sup:cusp}

We first separate the secondary triple-root cluster from the three simple
roots on the rest of the interval $I$. This localization supplies the
complete cycle count when the cluster is unfolded.

\subsection{Location of all roots near the secondary triple root}

\begin{lemma}[All roots near the secondary triple root]
\label{corner:all-roots}
Use the standard unfolding.
Fix a sufficiently small interval $J=(z_-,z_+)\ni1/3$.
There are a neighbourhood $\mathcal U$ of
$(-1/27,2/(9B))$, with compact closure in $\{\alpha<0,d>0\}$,
and $\mu_0,\sigma_0>0$ such that the following statements hold for every
\[
 (\alpha,d)\in\overline{\mathcal U},\qquad
 0<\mu<\mu_0,\qquad 0<\sigma<\sigma_0.
\]
Then $P''$ has exactly two simple zeros $\zeta_1<\zeta_2$ in the outer region,
with
\begin{equation}\label{corner:independent-curvature-roots}
 \sqrt{\zeta_1}=\mu\sigma\{\sqrt{Bd/2}+O(\mu+\sigma)\},\qquad
 \sqrt{\zeta_2}=\mu\{3/(8C)+O(\mu+\sigma^2)\}.
\end{equation}
The interval $\mathcal J=\{\mu^2\sigma^2z^2:z\in J\}$ contains
$\zeta_1$ and lies strictly between the second threshold and $\zeta_2$.
Outside $\mathcal J$ the displacement has exactly three roots,
all simple, with uniform asymptotics
\begin{equation}\label{corner:cusp-companions}
\begin{aligned}
 q_{N,1}&=-\frac{\mu^2\sigma}{C}(1+o(1)),&
 q_{N,2}&=\mu^2\sigma^2(\alpha+o(1)),\\
 q_{P,\mathrm{out}}&=\frac{\mu^2}{C^2}(1+o(1)).
\end{aligned}
\end{equation}
Inside $\mathcal J$ there are between one and three distinct roots.
All limits are uniform as $\mu+\sigma\to0$, without a condition on
the ratio of the two scales.
\end{lemma}

\begin{proof}
Use the separation graph~\eqref{corner:curvature-division} with
$p=(a,b,\xi_0)$. Substitution of $r=\mu\sigma z_g$ gives
\[
 d=z_g^2\{H_\Delta(\mu\sigma z_g,\mu,p)
                   +\sigma z_g B_\Delta(\mu\sigma z_g,p)\}.
\]
At zero scales the positive solution is $z_g=\sqrt{Bd/2}$, and the
state derivative of the right side is $4z_g/B>0$.
Substitution of $r=\mu w_g$ instead gives
\[
 \sigma^2d=w_g^2\{H_\Delta(\mu w_g,\mu,p)+w_gB_\Delta(\mu w_g,p)\}.
\]
Its nonzero positive solution at zero scales is $w_g=3/(8C)$; the
state derivative there is $-3/(4BC)$. The parameter implicit function
theorem proves~\eqref{corner:independent-curvature-roots} uniformly on
a compact neighbourhood of the parameters. In the second equation the restriction
at $\mu=0$ depends on $\sigma$ only through $\sigma^2d$, giving the
stated remainder. Both roots are in the outer region, since
$\delta/\zeta_1=O(\mu)$ and $\delta/\zeta_2=O(\mu\sigma^2)$.
The graph with its unique positive critical point and the threshold-layer
estimate of~\cite[Lemma~\ref*{main-lem:event-layer}]{main2026} exhaust the
interval and exclude any other zero of $P''$. These two simple zeros and the threshold-layer signs give
\[
\begin{array}{c|ccccc}
q&(-q_0,0)&(0,a_\delta)&(a_\delta,\zeta_1)&(\zeta_1,\zeta_2)&(\zeta_2,q_0)\\ \hline
\sgn D_{qq}&-&+&-&+&-
\end{array}
\]
where $a_\delta=\Gamma_1^{-1}(\delta)$.

The threshold identities $Q(a_\delta)=\delta$ and
$a_\delta=\delta-A_1\delta^{3/2}+O(\delta^2)$ yield
\[
 D(a_\delta)-a=b\delta+A_1\delta^{3/2}+O(\delta^2),\qquad
 \frac{D(a_\delta)}{\mu^4\sigma^3}=\alpha+O(\sqrt\mu).
\]
Indeed the three normalized errors are respectively
$O(\mu)$, $O(\sqrt\mu)$ and $O(\mu^2\sigma)$.
Thus $D(0)<0$ and $D(a_\delta)<0$ uniformly. Strict convexity on $(0,a_\delta)$
excludes all roots between and at the thresholds.

Choose $J=(1/3-\varepsilon,1/3+\varepsilon)$ with $\varepsilon>0$
fixed and small. At the central parameters the zero-scale quotient is
$(z-1/3)^3$. Its endpoint values are $\pm\varepsilon^3$ and its
endpoint derivatives are $3\varepsilon^2$. First shrinking
$\mathcal U$, then the common scale rectangle, gives
\begin{equation}\label{corner:cusp-window-signs}
 \mathcal F(z_-)<0<\mathcal F(z_+),\qquad
 \mathcal F_z(z_\pm)>0.
\end{equation}
Write $q_\pm=\mu^2\sigma^2z_\pm^2$. The curvature locations and
$a_\delta\sim\delta$ give $a_\delta<q_-<\zeta_1<q_+<\zeta_2$.
In particular,
\[
 D(q_-)<0<D(q_+),\qquad
 D_q(q_\pm)=\frac{\mu^2\sigma}{2z_\pm}
                       \mathcal F_z(z_\pm)>0.
\]
On $(a_\delta,q_-)$ the derivative is decreasing and hence positive, so
$D<D(q_-)<0$. On $(q_+,\zeta_2)$ it is increasing and positive, so
$D>D(q_+)>0$, including $D(\zeta_2)>0$.

Lemma~\ref{corner:realization-charts} gives one simple root near each
of $x=\alpha$, $X=1$, and $X=1/C^2$ in its three charts. Their
nonzero limiting derivatives have common margins on
$\overline{\mathcal U}$, and their isolators are disjoint for small
$\sigma$. This proves~\eqref{corner:cusp-companions}.
The negative half-interval is strictly concave and therefore has at most
two zeros, so the two negative roots exhaust it. The outer positive
root lies beyond $\zeta_2$, since $1/C>3/(8C)$, and has
$D_q=\mu^2\{-1/(2C)+o(1)\}<0$. Strict concavity on $(\zeta_2,q_0)$,
starting at the positive value $D(\zeta_2)$, keeps $D$ positive before
this root and negative after it. This excludes every other root up
to the fixed right endpoint. The negative-side concavity likewise
covers its entire fixed tail.

Only $\mathcal J$ remains. Its endpoint values have opposite signs,
and its only zero of $P''$ is $\zeta_1$, with negative curvature before
and positive curvature after it. Thus $D_q$ has at most two zeros there,
and Rolle's theorem gives at most three roots. All comparisons above
use separate small bounds on $\mu$ and $\sigma$ on the same interval.
\end{proof}

\subsection{The cusp and its fold curves}

Lemma~\ref{corner:all-roots} places the secondary cluster where both passages
are active, with three simple roots on the rest of $I$. We apply the
preparation theorem to the joint quotient of
Lemma~\ref{corner:joint-quotient}.

\begin{theorem}[Cusp and cycle counts]\label{thm:actual-cusp}
For the standard unfolding, there are a neighbourhood
$\mathcal U\ni(-1/27,2/(9B))$ and $\mu_0,\sigma_0>0$
such that the following hold on $\mathcal U$ for every independent pair
$0<\mu<\mu_0$, $0<\sigma<\sigma_0$. There is a unique nearby
jointly smooth branch $(\alpha_c,d_c,z_c)(\mu,\sigma)$ of triple zeros
of $\mathcal F_{\mu,\sigma}$ through
$(-1/27,2/(9B),1/3)$ at zero scales.
The separation $\delta_{\rm cusp}=\mu^3\sigma^2d_c$ and
state $q_{\rm cusp}=\mu^2\sigma^2z_c^2$ give a limit cycle of multiplicity
three.

In the normalized variables $(\mu,\sigma,z,\alpha,d)$ there are jointly
smooth changes of state and parameters, with a nonzero smooth factor, which
reduce the local zero equation to
\begin{equation}\label{corner:cusp-preparation}
 y^3+uy+v=0.
\end{equation}
For each positive $\mu,\sigma$, $(u,v)$ are local coordinates on the slice
$m=\mu$, $b=-\mu^2\sigma$ of the standard unfolding. Throughout $\mathcal U$,
the prepared coordinates give the following counts in $I$:
four simple cycles when $4u^3+27v^2>0$ and six simple cycles when
$4u^3+27v^2<0$. Both regions are nonempty and open. On the fold arcs
$4u^3+27v^2=0$, $(u,v)\ne(0,0)$, there are five distinct cycles,
one double; at $(u,v)=(0,0)$ there are four distinct cycles, one triple.

In the rescaled parameters $(\alpha,d)$ the double-root curve is obtained
by solving
\begin{equation}\label{corner:actual-fold-curve}
 \mathcal F(t;\alpha_f,d_f)=0,\qquad
 \mathcal F_z(t;\alpha_f,d_f)=0
\end{equation}
for $(\alpha_f,d_f)$, with $t$ near $1/3$.
The curve has an ordinary cusp at $t=z_c$ and consists
of two ordinary fold arcs away from that point.
\end{theorem}

\begin{proof}
At $(\mu,\sigma,\alpha,d,z)=(0,0,-1/27,2/(9B),1/3)$,
the quotient is $(z-1/3)^3$, and
\[
 D_{(\alpha,d,z)}(\mathcal F,\mathcal F_z,\mathcal F_{zz})
 =\begin{pmatrix}1&B/2&0\\0&3B/2&0\\0&0&6\end{pmatrix},
 \qquad \det=9B.
\]
The parameter implicit function theorem gives the jointly smooth branch,
which differs from its limiting value by $O(\mu+\sigma)$ and has
$\mathcal F_{zzz}=6+O(\mu+\sigma)>0$.
Moreover
\[
 \frac{q_{\rm cusp}}{\delta_{\rm cusp}}
 =\frac{z_c^2}{\mu d_c}
 =\frac B{2\mu}\{1+O(\mu+\sigma)\}\longrightarrow\infty.
\]
Since the second threshold is asymptotic to $\delta_{\rm cusp}$, both
passages are active at the triple root and $D$ is smooth in $q$ there; the
state change $q=\mu^2\sigma^2z^2$ preserves multiplicity.

The Malgrange preparation theorem~\cite[Chapter~IV, Section~2, Theorem~2.1 and Note~(1), p.~95]{golubitsky1973}, with
both scales as parameters, followed by removal of the quadratic term gives
\eqref{corner:cusp-preparation}. At zero scales, with
$\widehat\alpha=\alpha+1/27$ and $\widehat d=d-2/(9B)$,
\[
\begin{aligned}
 u&=\tfrac32B\widehat d+O((|\widehat\alpha|+|\widehat d|)^2),\\
 v&=\widehat\alpha+\tfrac B2\widehat d
                         +O((|\widehat\alpha|+|\widehat d|)^2).
\end{aligned}
\]
At zero scales the parameter Jacobian is $-3B/2\ne0$, so it remains
nonzero on a common reduction.
Choose $J$ of Lemma~\ref{corner:all-roots} inside the preparation chart,
then reduce the parameters and the rectangle of scales so that every prepared
cubic root pulls back into $J$. The cubic discriminant gives one or
three simple roots on its two open sides, two distinct roots on a
nonzero fold, and a triple root at the origin. Both open sides occur
because the parameter change is a local diffeomorphism. Adding the three
simple roots of Lemma~\ref{corner:all-roots} gives all stated counts
on $I$.

For the fold curve, the equations
$(\mathcal F,\mathcal F_z)=0$ at $z=t$ have parameter Jacobian
\[
 D_{(\alpha,d)}(\mathcal F,\mathcal F_z)
 =\begin{pmatrix}1&3Bt/2\\0&3B/2\end{pmatrix}
 \quad\hbox{at zero scales}.
\]
The uniform implicit function theorem gives jointly smooth functions
$(\alpha_f,d_f)$ on a compact $t$ interval, with zero-scale curve
$\alpha_0(t)=-t^2+2t^3$, $d_0(t)=2t(2-3t)/(3B)$.
Differentiating the fold equations gives $\gamma_f'(z_c)=0$ for
$\gamma_f=(\alpha_f,d_f)$, since their parameter Jacobian is invertible
and $\mathcal F_z=\mathcal F_{zz}=0$ there. Smooth dependence on the
scales and the zero-scale curve give
\begin{equation}\label{corner:ordinary-cusp-determinant}
 \det\bigl(\gamma_f''(z_c),\gamma_f'''(z_c)\bigr)
       =\frac{48}{B}+O(\mu+\sigma)\ne0.
\end{equation}
Thus $\gamma_f$ has an ordinary cusp. Along it,
$\frac{d}{dt}\mathcal F_{zz}=\mathcal F_{zzz}>0$ at the cusp;
after a common reduction all other points are ordinary double roots.
\end{proof}

Figure~\ref{fig:cusp-local-diagram} shows the cycle counts on $I$ in the
prepared coordinates.
% Discriminant of the prepared cubic y^3+uy+v.
% Curve parameter: (u,v)=(-3t^2,2t^3).
% The homogeneous rescaling (u,v)=(eta^2 U,eta^3 V)
% puts this drawing inside an arbitrarily small prepared-control chart.
% Source and independent polynomial checks: scripts/check-secondary.py.
\begin{figure}[!htbp]
\centering
\begin{tikzpicture}[x=7.7cm,y=3.6cm,font=\small,
  boundary/.style={black,line width=.75pt},
  guide/.style={black!55,line width=.35pt},
  leader/.style={black,line width=.4pt,-{Stealth[length=1.4mm]}},
  region/.style={fill=white,inner sep=2pt,font=\normalsize}]
  \begin{scope}
    \clip (-1.25,-.50) rectangle (.42,.50);
    \fill[black!9]
      plot[domain=-.67:.67,variable=\t,samples=181]
        ({-3*(\t)^2},{2*(\t)^3}) -- cycle;
    \draw[boundary]
      plot[domain=-.67:.67,variable=\t,samples=181]
        ({-3*(\t)^2},{2*(\t)^3});
  \end{scope}
  \draw[guide,-{Stealth[length=1.5mm]}]
    (-1.25,0)--(.47,0) node[right,black] {$u$};
  \draw[guide,-{Stealth[length=1.5mm]}]
    (0,-.50)--(0,.55) node[above,black] {$v$};
  \node[region,fill=black!9] at (-.88,.015) {$6$};
  \node[region] at (.26,.25) {$4$};
  \node[region] at (-.72,.44) {$4$};
  \node[anchor=west,font=\scriptsize] (fold) at (-.49,.34) {fold: $5$};
  \draw[leader] (fold.west)--(-.633,.196);
  \fill (0,0) circle[radius=1.8pt];
  \node[anchor=west,font=\scriptsize] (cusp) at (.11,-.20) {cusp: $4$};
  \draw[leader] (cusp.west)--(.008,-.008);
\end{tikzpicture}
\caption{The cusp in the prepared coordinates of
$y^3+uy+v=0$. The boundary is $4u^3+27v^2=0$:
the shaded region has six simple cycles and the other region has four,
including the three simple cycles outside the local cluster.
An ordinary fold has five different cycles, one double; the cusp has four,
one triple. A sufficiently small weighted rescaling of $(u,v)$ places the
diagram inside the common local parameter neighbourhood.
The pullback of the discriminant to $(\alpha,d)$ is the local bifurcation
curve, including its cusp.}
\label{fig:cusp-local-diagram}
\end{figure}

